\documentclass{article}
\usepackage{amsmath,amssymb,amsthm}
\usepackage{algorithm}
\usepackage{algorithmic}
\usepackage{enumitem}
\usepackage{hyperref}
\usepackage{geometry}
\geometry{margin=1in}
\newtheorem{assumption}{Assumption}
\newtheorem{theorem}{Theorem}
\newtheorem{lemma}{Lemma}
\newcommand{\diam}{\mathrm{diam}}
\newcommand{\gap}{\mathcal{G}}
\newcommand{\E}{\mathbb{E}}

\begin{document}

\section*{Algorithm Name}
\textbf{Conditional Gradient Method (Frank–Wolfe) for Non‑Convex Smooth Optimization}

\section*{Mathematical Formulation}
Let $f:\mathbb{R}^d\rightarrow\mathbb{R}$ be a differentiable (possibly non‑convex) objective and let 
\[
\mathcal{C}\subset\mathbb{R}^d
\]
be a non‑empty compact convex set (the feasible domain).  
Denote by 
\[
s_k \;:=\; \arg\min_{s\in\mathcal{C}}\;\big\langle \nabla f(x_k),\,s\big\rangle
\tag{1}
\]
the \emph{linear minimization oracle} (LMO).  
Given a stepsize sequence $\{\gamma_k\}_{k\ge 0}\subset(0,1]$, the update reads
\[
x_{k+1}\;=\;(1-\gamma_k)x_k+\gamma_k s_k
\;=\;x_k+\gamma_k\big(s_k-x_k\big).
\tag{2}
\]

\section*{Pseudocode}
\begin{algorithm}[H]
\caption{Frank–Wolfe for non‑convex smooth $f$}
\begin{algorithmic}[1]
\REQUIRE Initial point $x_0\in\mathcal{C}$, stepsize rule $\{\gamma_k\}$.
\FOR{$k=0,1,2,\dots$}
    \STATE Compute gradient $g_k\gets \nabla f(x_k)$.
    \STATE \textbf{LMO:} $s_k \gets \arg\min_{s\in\mathcal{C}}\langle g_k,s\rangle$.
    \STATE Update $x_{k+1}\gets (1-\gamma_k)x_k+\gamma_k s_k$.
    \IF{termination criterion satisfied (e.g. $\gap(x_k)\le\varepsilon$ or $k\ge K_{\max}$)}
        \STATE \textbf{break}
    \ENDIF
\ENDFOR
\RETURN $x_{k+1}$
\end{algorithmic}
\end{algorithm}
\noindent
The \emph{Frank–Wolfe gap} (a stationarity measure) is defined as
\[
\gap(x)\;:=\;\max_{s\in\mathcal{C}}\langle \nabla f(x),\,x-s\rangle
\;=\;\langle \nabla f(x),\,x-s(x)\rangle,
\tag{3}
\]
where $s(x)$ denotes any solution of the LMO (1).

\section*{Dry Run Example}
Consider the univariate problem
\[
\min_{w\in[0,5]}\;f(w):=(w-3)^2,
\]
with $x_0:=w_0=0$ and a \emph{constant} stepsize $\gamma_k\equiv\gamma=0.1$.

\begin{enumerate}[label=\textbf{Iter \arabic*:},leftmargin=*]
\item \textbf{Iteration $k=0$}
\begin{align*}
g_0 &=\nabla f(w_0)=2(w_0-3)= -6,\\
\text{LMO: } s_0 &=\arg\min_{s\in[0,5]}\langle -6,s\rangle
                =\arg\min_{s\in[0,5]}(-6s)=5,\\
w_{1} &= (1-0.1)w_0+0.1\,s_0 =0.9\cdot0+0.1\cdot5=0.5,\\
f(w_1)&=(0.5-3)^2=6.25.
\end{align*}

\item \textbf{Iteration $k=1$}
\begin{align*}
g_1 &=2(w_1-3)=2(0.5-3)=-5,\\
s_1 &=\arg\min_{s\in[0,5]}(-5s)=5,\\
w_{2} &=0.9\cdot0.5+0.1\cdot5=0.95,\\
f(w_2)&=(0.95-3)^2\approx4.2025.
\end{align*}

\item \textbf{Iteration $k=2$}
\begin{align*}
g_2 &=2(w_2-3)=2(0.95-3)=-4.1,\\
s_2 &=5,\\
w_{3} &=0.9\cdot0.95+0.1\cdot5=1.355,\\
f(w_3)&=(1.355-3)^2\approx2.714.
\end{align*}
\end{enumerate}
The objective value decreases at every step, illustrating the algorithmic mechanics.

\newpage
\section*{Assumptions}
\begin{assumption}[Compact feasible set]\label{as:compact}
The domain $\mathcal{C}$ is convex, compact and has finite diameter
\[
D\;:=\;\diam(\mathcal{C})\;=\;\max_{x,y\in\mathcal{C}}\|x-y\|_2<\infty .
\tag{4}
\]
\end{assumption}
\begin{assumption}[Smoothness]\label{as:smooth}
$f$ is $L$‑smooth on an open set containing $\mathcal{C}$, i.e.
\[
\|\nabla f(x)-\nabla f(y)\|_2\;\le\;L\|x-y\|_2\qquad\forall x,y\in\mathcal{C}.
\tag{5}
\]
Equivalently,
\[
f(y)\;\le\;f(x)+\langle\nabla f(x),y-x\rangle+\frac{L}{2}\|y-x\|_2^2,
\qquad\forall x,y\in\mathcal{C}.
\tag{6}
\]
\end{assumption}
\begin{assumption}[Existence of a global minimum]\label{as:opt}
Denote $f^\star:=\inf_{x\in\mathcal{C}}f(x)$. Because $\mathcal{C}$ is compact and $f$ is continuous, $f^\star$ is attained.
\tag{7}
\end{assumption}

No convexity of $f$ is required; the analysis holds for arbitrary (possibly non‑convex) smooth objectives.

\section*{Main Convergence Theorem}
\begin{theorem}[Convergence of Frank–Wolfe for non‑convex smooth $f$]\label{thm:main}
Assume \ref{as:compact}–\ref{as:opt} and let the stepsizes be a constant $\gamma_k\equiv\gamma\in(0,1]$. Then for any integer $K\ge0$ the iterates $\{x_k\}$ generated by \eqref{2} satisfy
\[
\frac{1}{K+1}\sum_{k=0}^{K}\gap(x_k)
\;\le\;
\frac{f(x_0)-f^\star}{\gamma (K+1)}\;+\;\frac{L D^{2}}{2}\,\gamma .
\tag{8}
\]
Consequently, choosing $\displaystyle\gamma:=\sqrt{\frac{2\big(f(x_0)-f^\star\big)}{L D^{2}(K+1)}}\in(0,1]$ yields the rate
\[
\min_{0\le k\le K}\gap(x_k)\;\le\;
\frac{2D\sqrt{L\big(f(x_0)-f^\star\big)}}{\sqrt{K+1}} .
\tag{9}
\]
Thus the Frank–Wolfe gap converges to zero at the $O(1/\sqrt{K})$ rate even when $f$ is non‑convex.
\end{theorem}

\section*{Theoretical Properties}
\begin{itemize}
\item \textbf{Stability.} The bound \eqref{8} depends only on the diameter $D$ and the smoothness constant $L$, both intrinsic to the problem data. No curvature of $f$ is required.
\item \textbf{Hyper‑parameter sensitivity.} The constant stepsize $\gamma$ appears linearly in the second term of \eqref{8} and inversely in the first term. The optimal choice \eqref{9} balances the two contributions, giving the familiar $O(1/\sqrt{K})$ rate.
\item \textbf{Comparison to baselines.}
    \begin{itemize}
        \item \emph{Projected Gradient Descent (PGD)} with the same $L$‑smoothness achieves $O(1/K)$ in the convex case but requires a Euclidean projection, which can be expensive.
        \item \emph{Frank–Wolfe} avoids projections; the price is a slower $O(1/\sqrt{K})$ rate in the non‑convex regime, matching the best known guarantees for projection‑free methods under the same assumptions.
    \end{itemize}
\item \textbf{Special case – convex $f$.} If $f$ is convex, the Frank–Wolfe gap upper‑bounds the primal suboptimality, and the same derivation gives the classical $O(1/K)$ convergence rate by using the diminishing stepsize $\gamma_k=2/(k+2)$.
\end{itemize}

\section*{Discussion}
The theorem shows that even without convexity the conditional gradient scheme drives the Frank–Wolfe gap to zero, i.e. the iterates become \emph{approximately stationary} for the constrained problem. The proof hinges only on smoothness and compactness, making it applicable to a broad class of machine‑learning problems where projection is costly but a cheap linear oracle is available (e.g., trace‑norm balls, simplex, flow polytopes).

\newpage
\section*{Preliminaries (Definitions and Lemmas)}
\begin{lemma}[Smoothness inequality]\label{lem:smooth}
If $f$ is $L$‑smooth on $\mathcal{C}$, then for any $x,y\in\mathcal{C}$
\[
f(y)\le f(x)+\langle \nabla f(x),y-x\rangle+\frac{L}{2}\|y-x\|_2^2 .
\tag{10}
\]
\end{lemma}
\begin{proof}
Inequality \eqref{10} is precisely the definition of $L$‑smoothness given in \eqref{6}. ∎
\end{proof}

\begin{lemma}[Bound on the displacement]\label{lem:diam}
For any $x\in\mathcal{C}$ and any $s\in\mathcal{C}$,
\[
\|s-x\|_2\le D .
\tag{11}
\]
\end{lemma}
\begin{proof}
Directly from the definition of the diameter \eqref{4}. ∎
\end{proof}

\section*{Main Proof (Step‑by‑Step Derivation)}
We prove Theorem \ref{thm:main}. Throughout the proof we keep the notation $g_k:=\nabla f(x_k)$ and $s_k$ as in \eqref{1}.

\begin{proof}
\textbf{[Step 1]} Apply Lemma \ref{lem:smooth} with $x=x_k$ and $y=x_{k+1}=x_k+\gamma (s_k-x_k)$:
\[
f(x_{k+1})\le f(x_k)+\langle g_k,\,\gamma(s_k-x_k)\rangle+\frac{L}{2}\|\gamma(s_k-x_k)\|_2^2 .
\tag{12}
\]

\textbf{[Step 2]} Pull out the scalar $\gamma$:
\[
f(x_{k+1})\le f(x_k)+\gamma\langle g_k, s_k-x_k\rangle+\frac{L\gamma^{2}}{2}\|s_k-x_k\|_2^{2}.
\tag{13}
\]

\textbf{[Step 3]} By definition of the LMO, $s_k$ minimizes the linear form $\langle g_k,s\rangle$ over $\mathcal{C}$, hence
\[
\langle g_k, s_k-x_k\rangle
= -\max_{s\in\mathcal{C}}\langle g_k, x_k-s\rangle
= -\gap(x_k) .
\tag{14}
\]

\textbf{[Step 4]} Substitute \eqref{14} into \eqref{13}:
\[
f(x_{k+1})\le f(x_k)-\gamma\,\gap(x_k)+\frac{L\gamma^{2}}{2}\|s_k-x_k\|_2^{2}.
\tag{15}
\]

\textbf{[Step 5]} Use Lemma \ref{lem:diam} to bound the squared displacement:
\[
\|s_k-x_k\|_2^{2}\le D^{2}.
\tag{16}
\]

\textbf{[Step 6]} Insert \eqref{16} into \eqref{15}:
\[
f(x_{k+1})\le f(x_k)-\gamma\,\gap(x_k)+\frac{L D^{2}}{2}\,\gamma^{2}.
\tag{17}
\]

\textbf{[Step 7]} Rearrange \eqref{17} to isolate the gap term:
\[
\gamma\,\gap(x_k)\le f(x_k)-f(x_{k+1})+\frac{L D^{2}}{2}\,\gamma^{2}.
\tag{18}
\]

\textbf{[Step 8]} Sum inequality \eqref{18} from $k=0$ to $K$:
\[
\gamma\sum_{k=0}^{K}\gap(x_k)
\le \sum_{k=0}^{K}\bigl(f(x_k)-f(x_{k+1})\bigr)+\frac{L D^{2}}{2}\,\gamma^{2}(K+1).
\tag{19}
\]

\textbf{[Step 9]} The telescoping sum collapses:
\[
\sum_{k=0}^{K}\bigl(f(x_k)-f(x_{k+1})\bigr)=f(x_0)-f(x_{K+1})\le f(x_0)-f^{\star},
\tag{20}
\]
where the last inequality follows from \eqref{7}.

\textbf{[Step 10]} Plug \eqref{20} into \eqref{19} and divide by $\gamma(K+1)$:
\[
\frac{1}{K+1}\sum_{k=0}^{K}\gap(x_k)
\le \frac{f(x_0)-f^{\star}}{\gamma(K+1)}+\frac{L D^{2}}{2}\,\gamma .
\tag{21}
\]
Equation \eqref{21} is precisely the bound \eqref{8} stated in the theorem. ∎

\section*{Rate Derivation (With Constants)}
From \eqref{21} we obtain a bound on the \emph{minimum} gap by the usual averaging argument:
\[
\min_{0\le k\le K}\gap(x_k)\;\le\;\frac{1}{K+1}\sum_{k=0}^{K}\gap(x_k).
\tag{22}
\]
Let $\Delta_0:=f(x_0)-f^{\star}>0$.  
Define the scalar
\[
\gamma^{\star}:=\sqrt{\frac{2\Delta_0}{L D^{2}(K+1)}} .
\tag{23}
\]
Since $0<\gamma^{\star}\le1$ for all $K$ large enough, we may set the constant stepsize $\gamma=\gamma^{\star}$ in \eqref{21}. Substituting yields
\[
\frac{1}{K+1}\sum_{k=0}^{K}\gap(x_k)
\le
\frac{\Delta_0}{\gamma^{\star}(K+1)}+\frac{L D^{2}}{2}\,\gamma^{\star}
=
\frac{\Delta_0}{K+1}\sqrt{\frac{L D^{2}(K+1)}{2\Delta_0}}
+\frac{L D^{2}}{2}\sqrt{\frac{2\Delta_0}{L D^{2}(K+1)}} .
\tag{24}
\]
Both terms simplify to the same expression, giving
\[
\frac{1}{K+1}\sum_{k=0}^{K}\gap(x_k)
\le
\frac{2 D\sqrt{L\Delta_0}}{\sqrt{K+1}} .
\tag{25}
\]
Using \eqref{22} we finally obtain \eqref{9}:
\[
\min_{0\le k\le K}\gap(x_k)\le
\frac{2 D\sqrt{L\big(f(x_0)-f^{\star}\big)}}{\sqrt{K+1}} .
\tag{26}
\]
Thus the Frank–Wolfe gap decays at the $O(1/\sqrt{K})$ rate.

\section*{Special Cases}
\begin{enumerate}[label=\arabic*.]
\item \textbf{Convex $f$}. If $f$ is convex, the gap satisfies $\gap(x_k)\ge f(x_k)-f^{\star}$, so \eqref{8} directly yields the classic $O(1/K)$ rate when the diminishing stepsize $\gamma_k=2/(k+2)$ is used (see e.g., Jaggi, 2013).
\item \textbf{Polyhedral $\mathcal{C}$}. When $\mathcal{C}$ is a polytope, the LMO can be solved in linear time, and the diameter $D$ is often known analytically, allowing the practitioner to set the optimal constant stepsize \eqref{23} without trial‑and‑error.
\item \textbf{Strongly smooth but non‑convex}. If $f$ satisfies the Polyak–Łojasiewicz (PL) condition on $\mathcal{C}$, one can combine \eqref{8} with the PL inequality to obtain a linear convergence rate; however, such a condition is beyond the scope of the present analysis.
\end{enumerate}

\end{document}